\begin{table}[htbp]
\centering
\small
\caption{Problema G (google/gemini-2.5-flash-lite): Fase E, acurácia do LLM contra o rótulo real por número de exemplos.}
\label{tab:bloco3_G_faseE_nshot_flashlite}
\begin{tabular}{lllrrrrr}
\toprule
Subproblema & Classes & Atributos & n = 0 & n = 4 & n = 10 & n = 20 & n = 40 \\
\midrule
G1 (normalizado) & Mulher / Homem & peso/altura & 31,1 (20,7) & 19,6 (24,1) & 21,9 (23,4) & 14,4 (25,7) & 10,4 (26,9) \\
G1 (normalizado) & Mulher / Homem & x1/x2 & 43,3 (17,0) & 68,9 (9,3) & 70,0 (9,0) & 61,5 (11,6) & 59,6 (12,1) \\
G1 (normalizado) & A / B & peso/altura & 30,0 (21,0) & 23,3 (23,0) & 18,1 (24,6) & 24,8 (22,6) & 24,1 (22,8) \\
G1 (normalizado) & A / B & x1/x2 & 15,6 (25,3) & 17,8 (24,7) & 16,3 (25,1) & 12,2 (26,3) & 12,6 (26,2) \\
\textbf{G2 (kg e m)} & \textbf{Mulher / Homem} & \textbf{peso\_kg/altura\_m} & \textbf{86,7 (4,0)} & \textbf{78,5 (6,4)} & \textbf{88,5 (3,4)} & \textbf{87,8 (3,7)} & \textbf{87,4 (3,8)} \\
G2 (kg e m) & Mulher / Homem & x1/x2 & 56,7 (13,0) & 73,3 (8,0) & 70,0 (9,0) & 74,1 (7,8) & 74,8 (7,6) \\
G2 (kg e m) & A / B & peso\_kg/altura\_m & 34,4 (19,7) & 33,0 (20,1) & 23,3 (23,0) & 18,9 (24,3) & 20,7 (23,8) \\
\bottomrule
\multicolumn{8}{p{0.95\linewidth}}{\footnotesize Fase E nas 30 amostras de teste de cada semente; o especialista é a melhor métrica da Fase A (eleita por acurácia vs real). Acurácia do LLM contra o rótulo real, erros em 30 entre parênteses; médias das sementes e sorteios. Em negrito: G2 com nomes semânticos.} \\
\end{tabular}
\end{table}
